\documentclass[aspectratio=169,11pt]{beamer}
% ------------------------------------------------
% THEME
% ------------------------------------------------
\usetheme{Madrid}
\usecolortheme{seahorse}

% ------------------------------------------------
% PACKAGES
% ------------------------------------------------

\usepackage[utf8]{inputenc}
\usepackage[T1]{fontenc}
\usepackage{lmodern}
\usepackage{amsmath,amssymb}
\usepackage{booktabs}
\usepackage{graphicx}
\usepackage{hyperref}

% ------------------------------------------------
% PRESENTATION INFORMATION
% ------------------------------------------------

\title[AI Cybercrime Forensics]
{Forensics for Cybercrimes Committed Using Artificial Intelligence}

\subtitle
{PhD Research Presentation}

\author[Peter M. Ngugi]
{Peter M. Ngugi}

\institute[Zetech University]
{
	School of Computing and Informatics\\
	Zetech University\\
	PhD in Computer Science
}

\date
{2026}

% ------------------------------------------------
% FOOTER
% ------------------------------------------------

\setbeamertemplate{blocks}[rounded][shadow=true]
\setbeamertemplate{navigation symbols}{}

\setbeamertemplate{footline}
{
	\leavevmode%
	\hbox{
		\begin{beamercolorbox}
			[wd=.33\paperwidth,ht=2.5ex,dp=1ex,left]
			{author in head/foot}
			\hspace{1em}Peter M. Ngugi
		\end{beamercolorbox}
		
		\begin{beamercolorbox}
			[wd=.34\paperwidth,ht=2.5ex,dp=1ex,center]
			{title in head/foot}
			PhD Computer Science
		\end{beamercolorbox}
		
		\begin{beamercolorbox}
			[wd=.33\paperwidth,ht=2.5ex,dp=1ex,right]
			{date in head/foot}
			\insertframenumber{} / \inserttotalframenumber
			\hspace{1em}
		\end{beamercolorbox}
	}
}

% ------------------------------------------------
% DOCUMENT
% ------------------------------------------------

\begin{document}
	
	% ------------------------------------------------
	% TITLE PAGE
	% ------------------------------------------------
	
	\begin{frame}
		\titlepage
	\end{frame}
	
	% ------------------------------------------------
	% OUTLINE
	% ------------------------------------------------
	
	\begin{frame}{Presentation Outline}
		
		\tableofcontents
		
	\end{frame}
	
	% ------------------------------------------------
	% SECTION 1
	% ------------------------------------------------
	
	\section{Introduction}
	
	\begin{frame}{Introduction}
		
		Artificial Intelligence is increasingly being incorporated
		into systems used for communication, automation,
		decision-making and cyber operations.
		
		\vspace{0.5cm}
		
		This creates a new forensic challenge:
		
		\begin{block}{Research Problem}
			How can digital investigators reliably identify,
			collect, preserve, analyse and present evidence
			when cybercrimes involve AI-generated or
			AI-assisted activity?
		\end{block}
		
	\end{frame}
	
	% ------------------------------------------------
	% SECTION 2
	% ------------------------------------------------
	
	\section{Problem Statement}
	
	\begin{frame}{Problem Statement}
		
		Traditional digital forensic approaches generally assume
		that evidence can be associated with identifiable
		devices, accounts, users or processes.
		
		\vspace{0.4cm}
		
		AI-assisted cybercrime introduces additional questions:
		
		\begin{itemize}
			
			\item Who actually performed the action?
			
			\item Was the action human-generated or AI-assisted?
			
			\item How can investigators establish provenance?
			
			\item Can AI-generated evidence be authenticated?
			
			\item How should attribution be established?
			
			\item How can such evidence be presented in court?
			
		\end{itemize}
		
	\end{frame}
	
	% ------------------------------------------------
	% SECTION 3
	% ------------------------------------------------
	
	\section{Research Objectives}
	
	\begin{frame}{Research Objectives}
		
		\textbf{General Objective}
		
		\begin{block}{Objective}
			To develop a digital forensic framework for investigating
			cybercrimes committed using Artificial Intelligence.
		\end{block}
		
		\vspace{0.4cm}
		
		\textbf{Specific Objectives}
		
		\begin{enumerate}
			
			\item Examine existing AI-assisted cybercrime techniques.
			
			\item Evaluate existing digital forensic approaches.
			
			\item Investigate methods for establishing evidence
			provenance and integrity.
			
			\item Develop a forensic framework for AI-assisted
			cybercrime investigations.
			
			\item Evaluate the proposed framework.
			
		\end{enumerate}
		
	\end{frame}
	
	% ------------------------------------------------
	% SECTION 4
	% ------------------------------------------------
	
	\section{Research Questions}
	
	\begin{frame}{Research Questions}
		
		\begin{enumerate}
			
			\item How is AI being used to facilitate cybercrime?
			
			\item What limitations exist in current digital forensic
			approaches?
			
			\item How can AI-generated evidence be authenticated?
			
			\item How can provenance and attribution be established?
			
			\item What forensic framework can support investigation
			and courtroom presentation?
			
		\end{enumerate}
		
	\end{frame}
	
	% ------------------------------------------------
	% SECTION 5
	% ------------------------------------------------
	
	\section{Literature Review}
	
	\begin{frame}{Literature Review}
		
		The literature will be examined across four major areas:
		
		\begin{columns}
			
			\column{0.48\textwidth}
			
			\begin{block}{Digital Forensics}
				\begin{itemize}
					\item Evidence acquisition
					\item Preservation
					\item Analysis
					\item Presentation
				\end{itemize}
			\end{block}
			
			\column{0.48\textwidth}
			
			\begin{block}{AI and Cybercrime}
				\begin{itemize}
					\item Generative AI
					\item Deepfakes
					\item Automated attacks
					\item AI-assisted fraud
				\end{itemize}
			\end{block}
			
		\end{columns}
		
	\end{frame}
	
	% ------------------------------------------------
	% SECTION 6
	% ------------------------------------------------
	
	\section{Theoretical Framework}
	
	\begin{frame}{Theoretical Framework}
		
		The study will examine relevant theories and formal
		models that can explain forensic evidence and attribution.
		
		\vspace{0.4cm}
		
		Potential theoretical foundations include:
		
		\begin{itemize}
			
			\item Digital Forensics Theory
			
			\item Information Security Theory
			
			\item Evidence and Authentication Principles
			
			\item Attribution Theory
			
			\item Computational and Formal Reasoning
			
		\end{itemize}
		
	\end{frame}
	
	% ------------------------------------------------
	% SECTION 7
	% ------------------------------------------------
	
	\section{Methodology}
	
	\begin{frame}{Research Methodology}
		
		\begin{block}{Research Approach}
			The study will employ an appropriate mixed-methods /
			computational research approach depending on the
			research phase and research questions.
		\end{block}
		
		\vspace{0.5cm}
		
		\begin{itemize}
			
			\item Literature analysis
			
			\item Digital forensic experimentation
			
			\item Dataset analysis
			
			\item AI-assisted cybercrime scenarios
			
			\item Framework development
			
			\item Experimental evaluation
			
		\end{itemize}
		
	\end{frame}
	
	% ------------------------------------------------
	% SECTION 8
	% ------------------------------------------------
	
	\section{Proposed Framework}
	
	\begin{frame}{Proposed Forensic Framework}
		
		\begin{center}
			
			\begin{block}{AI Cybercrime Forensic Process}
				
				Collection
				$\rightarrow$
				Preservation
				$\rightarrow$
				Authentication
				$\rightarrow$
				Analysis
				$\rightarrow$
				Attribution
				$\rightarrow$
				Presentation
				
			\end{block}
			
		\end{center}
		
	\end{frame}
	
	% ------------------------------------------------
	% SECTION 9
	% ------------------------------------------------
	
	\section{Contribution to Knowledge}
	
	\begin{frame}{Expected Contribution to Knowledge}
		
		The research aims to contribute:
		
		\begin{enumerate}
			
			\item A forensic framework for AI-assisted cybercrime.
			
			\item Methods for assessing provenance of AI-generated
			digital evidence.
			
			\item Approaches for authentication and integrity
			assessment.
			
			\item A structured approach to AI-assisted cybercrime
			attribution.
			
			\item Guidance for forensic investigators and
			legal practitioners.
			
		\end{enumerate}
		
	\end{frame}
	
	% ------------------------------------------------
	% SECTION 10
	% ------------------------------------------------
	
	\section{Conclusion}
	
	\begin{frame}{Conclusion}
		
		\begin{block}{Central Argument}
			AI changes not only how cybercrimes can be committed,
			but also how digital evidence must be interpreted,
			authenticated and attributed.
		\end{block}
		
		\vspace{0.5cm}
		
		The research therefore seeks to bridge the gap between:
		
		\[
		\boxed{
			\text{Artificial Intelligence}
			\quad+\quad
			\text{Cybercrime}
			\quad+\quad
			\text{Digital Forensics}
			\quad+\quad
			\text{Evidence}
		}
		\]
		
	\end{frame}
	
	% ------------------------------------------------
	% REFERENCES
	% ------------------------------------------------
	
	\section{References}
	
	\begin{frame}{References}
		
		\small
		
		\begin{thebibliography}{99}
			
			\bibitem{ref1}
			Author, A. (2025).
			\textit{Digital Forensics and Artificial Intelligence}.
			Journal of Cybersecurity Research.
			
			\bibitem{ref2}
			Author, B. (2025).
			\textit{Artificial Intelligence and Cybercrime}.
			International Journal of Digital Security.
			
		\end{thebibliography}
		
	\end{frame}
	
	% ------------------------------------------------
	% END
	% ------------------------------------------------
	
	\begin{frame}
		
		\begin{center}
			
			\Huge
			\textbf{Thank You}
			
			\vspace{0.8cm}
			
			\Large
			Questions and Discussion
			
			\vspace{1cm}
			
			\normalsize
			Peter M. Ngugi\\
			PhD in Computer Science\\
			Zetech University
			
		\end{center}
		
	\end{frame}
	
\end{document}